\chapter{허락으로 걷기: 엔도스랭크}

드디어 도전자를 소개할 차례예요.
도전자는 아주 간단한 생각에서 나왔어요.

\begin{center}
\begin{tcolorbox}[colback=white, colframe=allowcol, boxrule=0.8pt, arc=2mm, width=0.9\linewidth, halign=center]
\titlefont\bfseries 쪽지 돌리기 놀이를\\송금 지도가 아니라 허락 지도에서 하자.
\end{tcolorbox}
\end{center}

3장에서 허락에는 “나는 너를 믿어”라는 뜻이 담겨 있다고 했어요.
6장에서는 쪽지 놀이가 “믿을 만한 친구가 보낸 화살표를 더 무겁게” 쳐 준다는 것을 보았어요.
이 둘을 합치면, 믿을 만한 친구들에게서 많이 믿음을 받은 통이 높은 점수를 받게 돼요.
저는 이 방법에 \textbf{엔도스랭크}라는 이름을 붙였어요.
영어로 “보증하다”라는 뜻의 낱말 “엔도스”와 “순위”라는 뜻의 낱말 “랭크”를 합친 이름이에요.

\section{엔도스랭크의 지도 그리는 법}

\begin{enumerate}
\item 허락 줄만 골라요.
\item (스티커 종류, 허락한 통, 허락받은 통)마다 \emph{가장 마지막} 허락만 남겨요.
      허락을 0으로 거두었다면 그 화살표는 지워요.
\item 허락한 통에서 허락받은 통으로 화살표를 그리고, 허락한 양을 무게로 적어요.
      같은 두 통 사이에 여러 종류의 허락이 있으면 무게를 더해요.
\item 이 지도에서 쪽지 돌리기 놀이를 해요.
\end{enumerate}

송금 랭크와 다른 점은 딱 하나, 어떤 줄로 지도를 그리느냐예요.
쪽지 돌리기 놀이 자체는 똑같아요.
“선생님께” 카드도 똑같이 15장이에요.
이렇게 하면 두 방법을 공정하게 비교할 수 있어요.
점수가 다르게 나온다면 그건 놀이 방법 때문이 아니라 지도 때문이니까요.

\section{무엇이 달라질까요?}

\subsection*{1. 점수의 뜻이 달라져요}

송금 랭크 점수가 높은 통은 “스티커가 많이, 최근에 흘러 들어온 통”이에요.
가게처럼 바쁜 통이지요.
엔도스랭크 점수가 높은 통은 “믿을 만한 통들이 문을 열어 준 통”이에요.
두 점수는 서로 다른 것을 재요.
11장에서 이 차이를 숫자로 확인할 거예요.

\subsection*{2. 지도가 훨씬 작아져요}

3장에서 허락 줄은 송금 줄보다 훨씬 적다고 했지요.
같은 지갑들로 지도를 그려 보니 화살표 수가 이렇게 달랐어요.

\begin{center}
\small
\begin{tabular}{@{}lrr@{}}
\toprule
& 엔도스랭크(허락 지도) & 송금 랭크(송금 지도) \\
\midrule
점의 수 & 6,442 & 30,791 \\
화살표의 수 & 14,727 & 137,087 \\
\bottomrule
\end{tabular}
\end{center}

화살표가 9배쯤 적어요.
물을 흘려보낼 길이 적으면 계산도 빨라지겠지요.
12장에서 실제로 얼마나 빨랐는지 볼게요.

\subsection*{3. 아무 말도 할 수 없는 통이 생겨요}

허락을 한 번도 하거나 받지 않은 통은 허락 지도에 아예 나타나지 않아요.
그런 통에 대해서 엔도스랭크는 아무 말도 할 수 없어요.
점수가 0이 되지요.
제가 살펴본 지갑 5,521개 가운데 966개가 그랬어요.
송금 지도에 나타나지 않는 지갑은 381개였어요.
이것은 엔도스랭크의 분명한 약점이에요.

\begin{figure}[h]
\centering
\begin{tikzpicture}[scale=0.72, every node/.style={transform shape}]
  \begin{scope}
    \node[kid] (a) at (90:2.3) {지우};
    \node[kid] (b) at (162:2.3) {하준};
    \node[kid] (c) at (234:2.3) {서연};
    \node[kid] (d) at (306:2.3) {도윤};
    \node[kid] (e) at (18:2.3) {민서};
    \node[robot] (r) at (0,0) {로봇};
    \draw[send] (a) -- (b);
    \draw[send] (b) to[bend right=12] (e);
    \draw[send] (d) -- (c);
    \draw[send] (c) to[bend left=25] (a);
    \draw[send] (e) -- (d);
    \draw[send] (a) -- (r);
    \draw[send] (b) -- (r);
    \draw[send] (c) -- (r);
    \draw[send] (r) -- (d);
    \draw[send] (e) to[bend right=12] (b);
    \draw[send] (d) to[bend left=10] (b);
    \node[font=\bfseries, color=sendcol, align=center] at (0,-3.5) {송금 랭크의 지도\\[-1mm]\small 1등: 하준, 로봇};
  \end{scope}
  \begin{scope}[xshift=7.2cm]
    \node[kid] (a) at (90:2.3) {지우};
    \node[kid] (b) at (162:2.3) {하준};
    \node[kid] (c) at (234:2.3) {서연};
    \node[kid] (d) at (306:2.3) {도윤};
    \node[kid] (e) at (18:2.3) {민서};
    \node[robot, fill=allowcol!20] (r) at (0,0) {로봇};
    \draw[allow] (a) -- (r);
    \draw[allow] (b) -- (r);
    \draw[allow] (c) -- (r);
    \draw[allow] (e) -- (r);
    \draw[allow] (d) -- (c);
    \node[font=\bfseries, color=allowcol, align=center] at (0,-3.5) {엔도스랭크의 지도\\[-1mm]\small 1등: 로봇};
  \end{scope}
\end{tikzpicture}
\caption{같은 교실이라도 어느 지도에서 쪽지를 돌리느냐에 따라 점수가 달라져요.}
\end{figure}

\section{왜 이런 생각을 했을까요?}

제가 읽은 블록체인 평판 연구는 대부분 송금 기록을 썼어요.
송금은 많고, 눈에 잘 띄고, “돈이 오갔다”는 분명한 사실이니까요.
그런데 허락 기록은 거의 쓰이지 않았어요.
허락은 스티커가 움직이지 않는 기록이라 그냥 지나치기 쉬웠던 거예요.
저는 그 조용한 기록에 오히려 믿음이 더 또렷하게 담겨 있을지도 모른다고 생각했어요.
그것이 이 논문 제목의 앞부분, “ERC-20 허락 화살표를 페이지랭크에 넣기”의 뜻이에요.

\begin{deeper}[진짜 정의]
논문에서 엔도스랭크의 지도는 연구 기간 끝(또는 동결일)에 살아 있는 최신 허락으로 만들어요.
(토큰 $t$, 주인 $o$, 사용자 $s$)마다 마지막 승인 사건의 값을 $a^{t}_{o\to s}$라 하면,
화살표 $o \to s$의 무게는 $\sum_t a^{t}_{o\to s}$예요(값이 0인 것은 빼요).
놀이 규칙은 송금 랭크와 같아요: $d=0.85$, 고르게 나누는 되돌림, 같은 멈춤 조건.
위의 점과 화살표 수는 매칭 코호트 5,521개 지갑으로 그린 지도의 크기예요.
\end{deeper}

\begin{learned}
\begin{itemize}
\item 엔도스랭크는 허락 지도에서 하는 쪽지 돌리기 놀이예요. 이 논문이 제안한 방법이에요.
\item 송금 랭크와는 지도만 다르고 놀이 규칙은 같아서 공정하게 비교할 수 있어요.
\item 엔도스랭크의 지도는 화살표가 9배쯤 적어요.
\item 허락 기록이 없는 지갑에 대해서는 엔도스랭크가 아무 말도 할 수 없어요.
\end{itemize}
\end{learned}
